% !TeX TXS-program:compile = txs:///arara
% arara: lualatex: {shell: yes, synctex: no, interaction: batchmode}
% arara: lualatex: {shell: yes, synctex: no, interaction: batchmode} if found('log', '(undefined references|Please rerun|Rerun to get)')

\documentclass[english,11pt,a4paper]{article}
%\usepackage[utf8x]{inputenc}
%\usepackage[T1]{fontenc}
%\usepackage{DejaVuSerif}
%\usepackage[scale=1.1]{inconsolata}
\usepackage{customenvs}
\usepackage{xspace}
\usepackage{tabularx}
\usepackage{soul}
\usepackage{codehigh}
%\usepackage{fontawesome5}
\usepackage{lipsum}
\usepackage{fancyvrb}
\usepackage{fancyhdr}
\fancyhf{}
\renewcommand{\headrulewidth}{0pt}
\lfoot{\sffamily\small [customenvs]}
\cfoot{\sffamily\small - \thepage{} -}
\rfoot{\hyperlink{matoc}{\small\faArrowAltCircleUp[regular]}}
\usepackage{hologo}
\providecommand\tikzlogo{Ti\textit{k}Z}
\providecommand\TeXLive{\TeX{}Live\xspace}
\providecommand\PSTricks{\textsf{PSTricks}\xspace}
\let\pstricks\PSTricks
\let\TikZ\tikzlogo

\usepackage{hyperref}
\urlstyle{same}
\hypersetup{pdfborder=0 0 0}
\usepackage[margin=1.5cm]{geometry}
\setlength{\parindent}{0pt}

\def\TPversion{0.41f}
\def\TPdate{24/09/2025}
\usepackage{tcolorbox}
\sethlcolor{lightgray!25}
\NewDocumentCommand\MontreCode{ m }{%
	\hl{\vphantom{\texttt{pf}}\texttt{#1}}%
}
\usepackage[auto]{contour}

\usepackage{babel}

\begin{document}

\pagestyle{fancy}

\thispagestyle{empty}

\begin{center}
	\begin{minipage}{0.75\linewidth}
	\begin{tcolorbox}[colframe=yellow,colback=yellow!15]
		\begin{center}
			\renewcommand\arraystretch{1.25}
			\begin{tabular}{c}
				{\Huge \texttt{customenvs [en]}}\\
				\\
				{\Large Some custom environments,} \\
				{\Large or small patches.} \\
			\end{tabular}
			\renewcommand\arraystretch{1}
			
			\medskip
			
			{\small \texttt{Version \TPversion{} -- \TPdate}}
		\end{center}
	\end{tcolorbox}
\end{minipage}
\end{center}

\begin{center}
\begin{minipage}{0.85\linewidth}
	\begin{tcolorbox}[colframe=teal,colback=teal!10,halign=center,fontupper=\footnotesize]
		\texttt{\url{https://github.com/cpierquet/latex-packages/tree/main/customenvs}}
		
		~
		
		\texttt{\url{https://forge.apps.education.fr/pierquetcedric/packages-latex}}
	\end{tcolorbox}
\end{minipage}
\end{center}

\vspace*{5mm}

%\hrule
%
%\vspace*{2mm}
%
%{\centering Le nom du package vient de \textsf{ENVironnemenTS ALTernatifs}, avec le suffixe \textsf{FRançais}.\par}
%
%\vspace*{2mm}
%
%\hrule
%
%\vspace*{1cm}

\hrule

\phantomsection

%\hypertarget{matoc}{}

\tableofcontents

\vspace*{5mm}

\hrule

\vfill

\section{History}

\verb|v0.41f|:~~~Enhancements + new pictos

\verb|v0.41e|:~~~Enhancements with subpackages + new environements (see \textsf{[fr]} doc)

\verb|v0.41d|:~~~Enhancements with subpackages

\verb|v0.41c|:~~~Enhancements with subpackages

\verb|v0.41b|:~~~Bugfix + pre-compatibility with \textsf{fa5/fa6/fa7} + new pictograms

\verb|v0.41a|:~~~Option \textsf{noinlinegraphicx} for compatibility with \textsf{MikTeX} + Bugfix

\verb|v0.40f|:~~~\textsf{WhatsApp} style for 'Chat'

\verb|v0.40e|:~~~\textsf{customenvs-icons} \texttt{v0.1.0}

\verb|v0.40d|:~~~Code enhancements (compatibility  with twemojis) + \textsf{customenvs-tikzpictos} \texttt{v0.1.4}

\verb|v0.40c|:~~~PictoClippy (\textsf{customenvs-tikzpictos} \texttt{v0.1.3}) + Lengths macros

\verb|v0.40b|:~~~PictoCalendar (\textsf{customenvs-tikzpictos} \texttt{v0.1.2}) + enhancements

\verb|v0.40a|:~~~PictoTraffic (\textsf{customenvs-tikzpictos} \texttt{v0.1.1}) + enhancements

\verb|v0.3.7|:~~~Auxiliary package \textsf{customenvs-tikzpictos} for pictograms

\verb|v0.3.6|:~~~Picto \textit{bullseye+arrow}

\verb|v0.3.5|:~~~Bugfix + pre-compatibility with \textsf{fa5/fa6}

\verb|v0.3.4|:~~~Pictoskill

\verb|v0.3.3|:~~~Annotate image

\verb|v0.3.2|:~~~Alt version of title banner

\verb|v0.3.1|:~~~Box for MCQ

\verb|v0.3.0|:~~~Bugfix with \textsf{beamer}

\verb|v0.2.7|:~~~Key for mixing answers in MCQ

\verb|v0.2.6|:~~~Whell of skills, speedometer

\verb|v0.2.5|:~~~Bugfix with exercices (\textsf{[fr]} macro)

\verb|v0.2.4|:~~~Small box \textit{marker}

\verb|v0.2.3|:~~~Highway signs + sold banners (see \textsf{[fr]} doc)

\verb|v0.2.2|:~~~Flared arrow, with \TikZ

\verb|v0.2.1|:~~~Enhancements for \textit{stars skills} + AutoGrid for \TikZ\ (see \textsf{[fr]} doc)

\verb|v0.2.0|:~~~Skills with stars (\textsf{fontawesome5} or \TikZ)

\verb|v0.1.9|:~~~Title banner

\verb|v0.1.8|:~~~Score banner

\verb|v0.1.7|:~~~Small patch for \textsf{Vignette} macro (see \texttt{[fr]} documentation)

\verb|v0.1.6|:~~~Small patchs for \textsf{displayskip} + \textsf{pas-tableur} (see \texttt{[fr]} documentation)

\verb|v0.1.5|:~~~New macros for boxes with \textsf{tcolorbox} (see \texttt{[fr]} documentation)

\verb|v0.1.4|:~~~Create a SMS conversation

\verb|v0.1.3|:~~~Environment for exercise(s) (in french doc)

\verb|v0.1.2|:~~~Pencil of skills

\verb|v0.1.1|:~~~Skills table (only french for the moment...)

\verb|v0.1.0|:~~~Initial version

\vspace*{5mm}

\pagebreak

\section{The package customenvs}

\subsection{Idea}

The idea is to propose some classics environments with customizations (some are, for the moment, only in french):

\begin{itemize}
	\item write in \textit{multicols}, with spacings enhancements;
	\item present answers for a \textit{MCQ};
	\item create a list with \textit{choosen items} (randomly or by numbers);
	\item present a skill table.
\end{itemize}

\smallskip

The globa idea is ti propose \textit{user-friendly} environments, with explicit customizations, without using verbose syntax; but there's other solutions, using for example \MontreCode{\textbackslash vspace} ou \MontreCode{\textbackslash  setlength} or \MontreCode{spacingtricks} package.

\subsection{Loading}

The package loads within the preamble with \MontreCode{\textbackslash usepackage\{customenvs\}}.

Loaded packages are:

\begin{itemize}
	\item \MontreCode{xstring}, \MontreCode{simplekv}, \MontreCode{listofitems}, \MontreCode{randomlist} and \MontreCode{xintexpr};
	\item \MontreCode{enumitem};
	\item \MontreCode{multicol};
	\item \MontreCode{tabularray};
	\item \MontreCode{fontawesome};
\end{itemize}

Due to limitations, \MontreCode{enumitem}/\MontreCode{multicol}/\MontreCode{tabularray}/\MontreCode{fontawesome5/6}/\MontreCode{inlinegraphicx} can be \textit{un}loaded by \MontreCode{customenvs} (user must load them manually) via options:

\begin{itemize}
	\item \MontreCode{$\mathtt{\langle}$beamer$\mathtt{\rangle}$} for using with \textsf{beamer};
	\item \MontreCode{$\mathtt{\langle}$noenum$\mathtt{\rangle}$};
	\item \MontreCode{$\mathtt{\langle}$nomulticol$\mathtt{\rangle}$};
	\item \MontreCode{$\mathtt{\langle}$notblr$\mathtt{\rangle}$};
	\item \MontreCode{$\mathtt{\langle}$noinlinegraphicx$\mathtt{\rangle}$};
	\item \MontreCode{$\mathtt{\langle}$nofa$\mathtt{\rangle}$};
	\item \MontreCode{$\mathtt{\langle}$fa6$\mathtt{\rangle}$};
	\item \MontreCode{$\mathtt{\langle}$fa7$\mathtt{\rangle}$}.
\end{itemize}

\begin{codehigh}[language=latex/latex3,style/main=teal!25,style/code=teal!25]
%with all packages
\usepackage{customenvs}

%with option to no load some packages
\usepackage[option(s)]{customenvs}
\end{codehigh}

\pagebreak

\subsection{Subpackage customenvs-tikzpictos (v0.20a)}

The package \MontreCode{customenvs-tikzpictos}, loaded within \MontreCode{customenvs} (but can be loaded independently), proposes small pictograms.

\begin{codehigh}[language=latex/latex3,style/main=teal!25,style/code=teal!25]
%\usepackage{customenvs-tikzpictos} %only if for standalone

\tkzpicto%
    [keys]
    <tikz options>
    {type=params}

%type= wifi/network/stars/speedo/bullseye/skills/pill/calendar
%params= nb/nblevels (except bullseye) or day/month (calendar)
%key height= len / auto (without depth) / dauto (with depth)
\end{codehigh}

\hfill\begin{tblr}{hlines,vlines,cells={font=\sffamily\Large}}
	Wifi & \pictowifi[colors=purple/yellow]{0}|\pictowifi[colors=purple/yellow]{1}|\pictowifi[colors=purple/yellow]{2}|\pictowifi[colors=purple/yellow]{3}|\pictowifi[colors=purple/yellow]{4}\\
	Wifi (bars) & \pictowifi[colors=purple/yellow,style=bar]{0}|\pictowifi[colors=purple/yellow,style=bar]{1}|\pictowifi[colors=purple/yellow,style=bar]{2}|\pictowifi[colors=purple/yellow,style=bar]{3}|\pictowifi[colors=purple/yellow,style=bar]{4}\\
	Network & \pictonetwork[colors=teal/orange]{0}|\pictonetwork[colors=teal/orange]{1}|\pictonetwork[colors=teal/orange]{2}|\pictonetwork[colors=teal/orange]{3}|\pictonetwork[colors=teal/orange]{4}\\
	Stars & \pictostars[colors=brown/cyan]{0}|\pictostars[colors=brown/cyan]{0.5}|\pictostars[colors=brown/cyan]{3}\\
	Speedometer & \pictospeedometer[colors=violet/olive]{}|\pictospeedometer[colors=violet/olive]{0}|\pictospeedometer[colors=violet/olive]{0.5}|\pictospeedometer[colors=violet/olive]{5.5}|\pictospeedometer[colors=violet/olive]{6} \\
	BullsEye & \pictobullseye[colors]|\pictobullseye[colors]<rotate=90>|\pictobullseye[colors]<rotate=180>|\pictobullseye[colors]<rotate=-90>\\
	Battery & \pictobattery[barcolor=auto]{}|\pictobattery[barcolor=auto]{1}|\pictobattery[barcolor=auto]{2}|\pictobattery[barcolor=auto]{3}\\
	Battery (flip) & \pictobattery[flip,barcolor=auto]{}|\pictobattery[flip,barcolor=auto]{1}|\pictobattery[flip,barcolor=auto]{2}|\pictobattery[flip,barcolor=auto]{3}\\
	Skills & \pictoskills[colors={red}]{}|\pictoskills[colors={red}]{1}|\pictoskills[colors={red}]{2}|\pictoskills[colors={red}]{3}\\
	Pill & \pictopill[colors=blue/red]{}|\pictopill[colors=blue/red]{1}|\pictopill[colors=blue/red]{3}\\
	TrafficLight & \pictotraffic{}|\pictotraffic{1}|\pictotraffic{2}|\pictotraffic{3}|\pictotraffic{123}|\pictotraffic{13}\\
	MiniCalendar & \pictocalendar{15/aug.}|\pictocalendar[color=red]{30/dec.}|\pictocalendar[color=blue]{14/feb.}\\
	Clippy & \pictoclippy|\pictoclippy[style=angry]|\pictoclippy[style=happangry]|\pictoclippy[style=confused]\\
	Dball & \pictodball{}|\pictodball{1}|\pictodball{2}|\pictodball{6}|\pictodball{7}\\
\end{tblr}\hfill\null

\pagebreak

\subsection{Subpackage customenvs-icons (v0.1.1)}

\MontreCode{customenvs} loads, for \textit{small} icons, \MontreCode{customenvs-icons} package.

The idea is to propose small icons, independently of \MontreCode{customenvs}.

\begin{codehigh}[language=latex/latex3,style/main=teal!25,style/code=teal!25]
\usepackage{customenvs-icons}       %only if for standalon

\ceicon%
  [%
    educ=TF,                        %boolean style 'educ'
    design=TF                       %boolean style 'design'
    height=...,                     %(d)auto / height / height+depth
    (d)strut=...                    %box choices (for dim calc)
  ]%
  <includegraphics options>%
  {nom}
\end{codehigh}

\begin{demohigh}[language=latex/latex3,style/main=teal!25,style/code=teal!25]
{\Huge X\ceicon{brush}\ceicon[height=auto]{brush}X}
\end{demohigh}

\begin{demohigh}[language=latex/latex3,style/main=teal!25,style/code=teal!25]
\ceicon[height=2cm,design]{browser-html}
\end{demohigh}

\begin{demohigh}[language=latex/latex3,style/main=teal!25,style/code=teal!25]
\ceicon[height=1in,mathsymb]{pyramid}
\end{demohigh}

\newpage

\section{Answers for a MCQ}

\subsection{Idea}

The idea is to propose an environment to present answers for a MQC with \MontreCode{tabularray} (and not \MontreCode{multicols}).

\smallskip

It's possible to use 2, 3 or 4 answers (and with 4 answers it's possible to use 2 columns.)

\begin{codehigh}[language=latex/latex3,style/main=teal!25,style/code=teal!25]
\AnswersMCQ[options]{list of answers}<tblr options>
\end{codehigh}

The avalailable \MontreCode{options} are:

\begin{itemize}
	\item \MontreCode{Width}: \MontreCode{0.99\textbackslash linewidth} by default;
	\item \MontreCode{Lines}: \MontreCode{false} by default;
	\item \MontreCode{SpaceCR} for Columns/Rows spacing, within \MontreCode{col/row} or \MontreCode{global}: \MontreCode{6pt/2pt} by default;
	\item \MontreCode{NumCols}, 2 or 4: \MontreCode{4} by default;
	\item \MontreCode{Labels} for the labels: \MontreCode{a.} by default;
	\begin{itemize}
		\item with \MontreCode{box} to use a \textit{Box};
		\item with \MontreCode{a} to \textit{enumerate} \MontreCode{a b c d};
		\item with \MontreCode{A} to \textit{enumerate} \MontreCode{A B C D};
		\item with \MontreCode{1} to \textit{enumerate} \MontreCode{1 2 3 4};
	\end{itemize}
	\item \MontreCode{FontLabels}: \MontreCode{\textbackslash bfseries} by default;
	\item \MontreCode{SpaceLabels}: \MontreCode{\textbackslash kern5pt} by default;
	\item \MontreCode{Shuffle}, for mixing answers: \MontreCode{false} by default;
	\item \MontreCode{Swap}, for ACBD instead of ABCD: \MontreCode{false} by default.
\end{itemize}

The list of answers must be given within \MontreCode{answA § answB § ...}.

Specific options for \MontreCode{tblr} are given between last optionnal argument, between \MontreCode{<...>}.

\subsection{Examples}

\begin{demohigh}[language=latex/latex3,style/main=teal!25,style/code=teal!25]
%default output
\AnswersMCQ{Answer A § Answer B § Answer C § Answer D}
\end{demohigh}

\begin{demohigh}[language=latex/latex3,style/main=teal!25,style/code=teal!25]
\AnswersMCQ[Lines]{Answer A § Answer B § Answer C § Answer D}

\AnswersMCQ[Lines,Shuffle]{Answer A1 § Answer B1 § Answer C1 § Answer D1}

\AnswersMCQ[Lines,Shuffle]{Answer A2 § Answer B2 § Answer C2 § Answer D2}
\end{demohigh}

\begin{demohigh}[language=latex/latex3,style/main=teal!25,style/code=teal!25]
\AnswersMCQ[Lines,Labels=(1.),SpaceLabels={~~~}]{Answer A § Answer B § Answer C}
\end{demohigh}

\begin{demohigh}[language=latex/latex3,style/main=teal!25,style/code=teal!25]
\AnswersMCQ[Labels={A.},FontLabels={\color{red}\bfseries}]%
    {Answer A § Answer B § Answer C § Answer D}
\end{demohigh}

\begin{demohigh}[language=latex/latex3,style/main=teal!25,style/code=teal!25]
\AnswersMCQ[Labels={1.},FontLabels={\color{red}\bfseries}]%
    {Answer A § Answer B § Answer C § Answer D}
\end{demohigh}

\begin{demohigh}[language=latex/latex3,style/main=teal!25,style/code=teal!25]
\AnswersMCQ[NumCols=2,Labels={A.},FontLabels={\color{red}\bfseries}]%
    {Answer A § Answer B § Answer C § Answer D}
\end{demohigh}

\begin{demohigh}[language=latex/latex3,style/main=teal!25,style/code=teal!25]
\AnswersMCQ[NumCols=2,Swap,Labels={A.},FontLabels={\color{red}\bfseries}]%
    {Answer A § Answer B § Answer C § Answer D}
\end{demohigh}

\begin{demohigh}[language=latex/latex3,style/main=teal!25,style/code=teal!25]
\AnswersMCQ[Lines,NumCols=2,SpaceCR=6pt/10pt,Labels=box]%
    {Answer A § Answer B § Answer C § Answer D}
\end{demohigh}

\begin{demohigh}[language=latex/latex3,style/main=teal!25,style/code=teal!25]
% checkedbox is \def\MCQanswersbox{\raisebox{-0.2ex}{\faSquare[regular]}}
\AnswersMCQ[Width=10cm,NumCols=2,Lines]%
    {$\displaystyle\frac1x$ § $1+\displaystyle\frac1x$ § $-2x^2+5$ § $-\infty$}
    <rows={1.5cm}>
\end{demohigh}

\pagebreak

\section{List with picked elements (random or not)}

\subsection{Global use}

The idea is to:

\begin{itemize}
	\item create a list of items, the base for choices;
	\item print the list with picked items.
\end{itemize}

\begin{codehigh}[language=latex/latex3,style/main=teal!25,style/code=teal!25]
\CreateItemsList{list}{macro}{listname}
\end{codehigh}

\begin{codehigh}[language=latex/latex3,style/main=teal!25,style/code=teal!25]
\ListItemsChoice[keys]{macro}{listname}(numbers)<enumitem options>!beamer options!
\end{codehigh}

The available \MontreCode{keys} are:

\begin{itemize}
	\item \MontreCode{Type}: \MontreCode{enum} or \MontreCode{item};
	\item \MontreCode{Random}: \MontreCode{false} by default.
\end{itemize}

The second argument, mandatory and between \MontreCode{\{...\}} is the macro for the list.

The third argument, mandatory and between \MontreCode{\{...\}} is the name of the list.

The fourth argument, mandatory and between \MontreCode{(...)} give:

\begin{itemize}
	\item the number of random items to display, with \MontreCode{Random=true};
	\item the numbers of picked itemps, within \MontreCode{num1,num2,...}.
\end{itemize}

The next argument, optional and between \MontreCode{<...>} gives specific options to \MontreCode{enumitem} environment.

The last argument, between \MontreCode{!..!} gives specific options to \MontreCode{enumitem} environment with \MontreCode{beamer}.

\medskip

Controls are done:

\begin{itemize}
	\item to verify that the list doesn't exist (for the creation);
	\item to verify that that the list still exist (for the display).
\end{itemize}

\subsection{Examples}

\begin{codehigh}[language=latex/latex3,style/main=teal!25,style/code=teal!25]
%creation of list ListItems, with macro \mylistofitems
\CreateItemsList%
    {Answer A,Answer B,Answer C,Answer D,Answer E,Answer F,Answer G,Answer H}%
    {\mylistofitems}{ListItems}
\end{codehigh}

\CreateItemsList%
{Answer A,Answer B,Answer C,Answer D,Answer E,Answer F,Answer G,Answer H}%
{\mylistofitems}{ListItems}

\begin{demohigh}[language=latex/latex3,style/main=teal!25,style/code=teal!25]
%items random
\ListItemsChoice[Random]{\mylistofitems}{ListItems}(5)
\end{demohigh}

\begin{demohigh}[language=latex/latex3,style/main=teal!25,style/code=teal!25]
%items picked 
\ListItemsChoice{\mylistofitems}{ListItems}(1,4,3,8,2)
\end{demohigh}

\begin{codehigh}[language=latex/latex3,style/main=teal!25,style/code=teal!25]
%creation of list ListItemsB, with macro \mylistofitemsb
\CreateItemsList%
    {{$\int_0^1 x^2 dx$},{$\int_0^1 x^3 dx$},{$\int_0^1 x^4 dx$},...}%
    {\mylistofitemsb}{ListItemsB}
\end{codehigh}

\CreateItemsList%
{{$\int_0^1 x^2 dx$},{$\int_0^1 x^3 dx$},{$\int_0^1 x^4 dx$},{$\int_0^1 x^5 dx$},{$\int_0^1 x^6 dx$},{$\int_0^1 x^7 dx$},{$\int_0^1 x^8 dx$}}%
{\mylistofitemsb}{ListItemsB}

\begin{codehigh}[language=latex/latex3,style/main=teal!25,style/code=teal!25]
%items picked 
\ListItemsChoice[Type=item]{\mylistofitemsb}{ListItemsB}(7,2,1,5,3)<label=$--$>
\end{codehigh}

\ListItemsChoice[Type=item]{\mylistofitemsb}{ListItemsB}(7,2,1,5,3)<label=$--$>

\pagebreak

\section{Pencil of skills}

\subsection{Global use}

The idea is to:

\begin{itemize}
	\item present of list of categories and skills;
	\item presented like a pencil.
\end{itemize}

The code (within CC-BY-SA 4.0 license) is adapted from:

\hfill{\footnotesize \url{https://tex.stackexchange.com/questions/504092/replicating-a-fancy-bordered-text-style-in-latex/504145#504145}}\hfill~

\begin{codehigh}[language=latex/latex3,style/main=teal!25,style/code=teal!25]
\PencilSkills[keys]<tikz options>{listofskills}
\end{codehigh}

The style is globally fixed, but there's some customization available.

\subsection{The macro}

Available \MontreCode{keys} are:

\begin{itemize}
	\item \MontreCode{FontCateg}: font for the categories;
	\item \MontreCode{FontBlock}: font for the skills;
	\item \MontreCode{Colors}: list of category's colors
	
	\MontreCode{BgCateg1/FgCateg1,BgCateg1/FgCateg1,...}
	
	(if \MontreCode{FgCateg1} est missing, \MontreCode{black} is used)
	\item \MontreCode{BlockWidth}: width of skill's block;
	\item \MontreCode{Scale}: global scale
	\item \MontreCode{BlackWhite}: boolean for B\&W.
\end{itemize}

The second argument, optional and between \MontreCode{<...>} gives specific options to \MontreCode{enumitem} environment.

\smallskip

The last argument, mandatory and between \MontreCode{(...)} give the list of categories/skills, within

\MontreCode{Categ1/ListSkills1,Categ2/ListSkills2,...}.

\subsection{Examples}

\begin{demohigh}[language=latex/latex3,style/main=teal!25,style/code=teal!25]
%default output
\PencilSkills{Search/Skill 1\\ Skill 2,Model/{Skill 1\\ Skill 2}}
\end{demohigh}

\begin{demohigh}[language=latex/latex3,style/main=teal!25,style/code=teal!25]
\PencilSkills[Scale=0.75]%
    {Search/Skill 1\\Skill 2,Model/{Skill 1\\Skill 2},%
    Represent/{Skill 1\\Skill 2},Calculate/{Skill 1\\Skill 2},%
    Reason/{Skill 1\\Skill 2},Communicate/{Skill 1\\Skill 2}}
\end{demohigh}

\begin{demohigh}[language=latex/latex3,style/main=teal!25,style/code=teal!25]
\PencilSkills[Scale=0.75,BlockWidth=3cm]<rotate=90>{
    Search/Skill 1\\Skill 2,Model/{Skill 1\\Skill 2}}
\hspace{1cm}
\PencilSkills[Scale=0.75,BlockWidth=3cm]<rotate=-90>{
    Search/Skill 1\\Skill 2,Model/{Skill 1\\Skill 2}}
\hspace{1cm}
\PencilSkills[Scale=0.75,BlockWidth=3cm,BlackWhite]<rotate=45>{
    Search/Skill 1\\Skill 2,Model/{Skill 1\\Skill 2}}
\end{demohigh}

\pagebreak

\section{Score banner}

\subsection{Global use}

The idea is to insert a score banner, with customization.

\begin{codehigh}[language=latex/latex3,style/main=teal!25,style/code=teal!25]
ScoreBanner[keys]{number}
\end{codehigh}

\begin{demohigh}[language=latex/latex3,style/main=teal!25,style/code=teal!25]
%default output
\ScoreBanner{}
\end{demohigh}

\subsection{The macro}

Available \MontreCode{keys} are:

\begin{itemize}
	\item \MontreCode{Height}: height of the banner (without the legend); \MontreCode{1} by default
	\item \MontreCode{Ratio}: ratio of boxes; \MontreCode{0.6} by default
	\item \MontreCode{Symbols}: labels; \MontreCode{A,B,C,D,E} by default
	\item \MontreCode{Legend}: legend (uppercase); \MontreCode{score} by default;
	\item \MontreCode{Font}: global font; \MontreCode{\textbackslash bfseries\textbackslash sffamily} by default
	\item \MontreCode{ShowLegend}: boolean for the legend;  \MontreCode{false} by default;
	\item \MontreCode{Colors}: colors for boxes;
	
	\hfill\MontreCode{colorNS1,colorNS2,colorNS3,colorNS4,colorNS5} by default;
	\item \MontreCode{ScaleSymbols}: scale H/V of labels;  \MontreCode{1.25,1.65} by default;
	\item \MontreCode{Colbg}: background color for select box;  \MontreCode{white} by default.
\end{itemize}

\smallskip

If the list of colors doesn't fill all the boxes, \MontreCode{lightgray} color is used.

\begin{demohigh}[language=latex/latex3,style/main=teal!25,style/code=teal!25]
\ScoreBanner[Legend=Geometry,Height=2]{4}
\end{demohigh}

\begin{demohigh}[language=latex/latex3,style/main=teal!25,style/code=teal!25,style/demo=yellow!25]
%bg of lower part is yellow!25
\def\lstcouleurs{colorNS1,colorNS2,colorNS3,colorNS4,colorNS5,purple}
\ScoreBanner%
  [ScaleSymbols={1.33,2},Height=3.25,ShowLegend=false,Ratio=0.75,
  Symbols={1,2,3,4,5,6},Colors=\lstcouleurs,
  Colbg=yellow!25]{1}
\end{demohigh}

\pagebreak

\section{SMS conversation}

\subsection{Global use}

The idea is to present a conversation of SMS.

\begin{codehigh}[language=latex/latex3,style/main=teal!25,style/code=teal!25]
\begin{ChatSMS}[keys]{name}
  \InSMS(*){time}{msg}
  \OutSMS*(*){time}{msg}
\end{ChatSMS}
\end{codehigh}

The style is globally fixed, but there's some customization available.

\subsection{The environment}

Available \MontreCode{keys} are:

\begin{itemize}
	\item \MontreCode{height}: height of the window (auto or specific); \MontreCode{auto} by default
	\item \MontreCode{width}: width of the window; \MontreCode{7cm} by default
	\item \MontreCode{margin}: margin (L or R) for the bubble \MontreCode{1.5cm} by default
	\item \MontreCode{color}: \textit{main} color (banner); \MontreCode{teal!75!cyan!75!white} by default;
	\item \MontreCode{colback}: color for background; \MontreCode{lightgray!5} by default
	\item \MontreCode{colorin}: color for incoming SMS; \MontreCode{lime!25} by default
	\item \MontreCode{colorout}: color for outcoming SMS; \MontreCode{teal!25} by default
	\item \MontreCode{writetxt}: text of sending zone; \MontreCode{Write} by default
	\item \MontreCode{fonttxt}: bubble's font; \MontreCode{\textbackslash normalfont} by default
	\item \MontreCode{avatar}: avatar of contact; \MontreCode{\textbackslash faAddressCard} by default
	\item \MontreCode{dispavatar}: boolean for displaying avatar near the bubbles; \MontreCode{false} by default
	\item \MontreCode{blackwhite}: boolean pour black\&white. \MontreCode{false} by default
\end{itemize}

The argument, mandatory and between \MontreCode{(...)} give the name of the contact.

\subsection{Macros for the bubbles}

Regarding the bubble creation commands, \MontreCode{\textbackslash InSMS} and \MontreCode{\textbackslash OutSMS}:

\begin{itemize}
	\item the \textit{starred} version does not display the \textit{checkmarks} of \textit{good reception};
	\item the first mandatory argument is the time to display;
	\item the second mandatory argument is the message to display (including multi-lines).
\end{itemize}

\subsection{Examples}

\begin{demohigh}[language=latex/latex3,style/main=teal!25,style/code=teal!25]
%with a personal image
\begin{ChatSMS}%
  [width=6cm,fonttxt=\sffamily,height=10cm,avatar=img/android,dispavatar]{CP}
  \InSMS{19:23}{Hi !}
  \OutSMS{19:23}{Hi !\\ How are you ?}
  \InSMS{19:25}{Just a problem with a math question\ldots}
  \OutSMS{19:26}{Wanna help ??}
  \InSMS{19:28}{Yes, I need to compute $\mathsf{\int_{0}^{1} x^2e^{-x}\,dx}$\ldots}
  \OutSMS*{19:30}{Take care !!}
\end{ChatSMS}
\end{demohigh}

\begin{demohigh}[language=latex/latex3,style/main=teal!25,style/code=teal!25]
\begin{ChatSMS}%
  [width=8cm,fonttxt=\sffamily,avatar=\faCanadianMapleLeaf,blackwhite]{CP}
  \InSMS{19:23}{Hi !}
  \OutSMS{19:23}{Hi !\\ How are you ?}
  \InSMS{19:25}{Just a problem with a math question\ldots}
  \OutSMS{19:26}{Wanna help ??}
  \InSMS{19:28}{Yes, I need to compute $\mathsf{\int_{0}^{1} x^2e^{-x}\,dx}$\ldots}
  \OutSMS*{19:30}{Take care !!}
\end{ChatSMS}
\end{demohigh}

\subsection{Style WhatsApp}

Un style type \textit{WhatsApp} est également disponible, avec un fonctionnement similaire à celui présenté précédemment.

Les clés disponibles sont:

\begin{itemize}
	\item \MontreCode{height}: \MontreCode{auto} by default
	\item \MontreCode{width}: \MontreCode{5cm} by default
	\item \MontreCode{bgcolor}: \MontreCode{lightgray!50} by default
	\item \MontreCode{receivecolor}: \MontreCode{greenwa!66!white} by default
	\item \MontreCode{sendcolor}: \MontreCode{white} by default
	\item \MontreCode{txtwrite}: \MontreCode{Message...} by default
	\item \MontreCode{fonttxt}: \MontreCode{sffamily} by default
	\item \MontreCode{avatar}: \MontreCode{\textbackslash faAddressCard} by default
	\item \MontreCode{showavatar}: \MontreCode{false} by default
	\item \MontreCode{bw}: \MontreCode{false} by default
	\item \MontreCode{txtwidth}: \MontreCode{0.55} by default.
\end{itemize}

\begin{demohigh}[language=latex/latex3,style/main=teal!25,style/code=teal!25]
\begin{EnvChatWA}{C. Pierquet}
	\WaDate{monday}
	\WaRec{19:23}{Hi!}
	\WaDate{today}
	\WaSend*{08:33}{Hi!!}
	\WaSend*{08:35}{You're using WhApp so:-)}
	\WaRec{10:15}{Yep!!}
	\WaRec{10:18}{Could you help me?}
	\WaRec{10:18}{I've got pb with tables in \LaTeX\ldots}
\end{EnvChatWA}
\end{demohigh}

\begin{demohigh}[language=latex/latex3,style/main=teal!25,style/code=teal!25]
\begin{EnvChatWA}[bw,showavatar,avatar=Image/avatar]{C. Pierquet}
	\WaDate{monday}
	\WaRec{19:23}{Hi!}
	\WaDate{today}
	\WaSend*{08:33}{Hi!!}
	\WaSend*{08:35}{You're using WhApp so:-)}
	\WaRec{10:15}{Yep!!}
	\WaRec{10:18}{Could you help me?}
	\WaRec{10:18}{I've got pb with tables in \LaTeX\ldots}
\end{EnvChatWA}
\end{demohigh}

\pagebreak

\section{Title banner}

\subsection{Global usage}

The idea is to propose a banner, made with \TikZ, to present for example a title.

The global style is fixed, but few customization are possible.

\begin{codehigh}[language=latex/latex3,style/main=teal!25,style/code=teal!25]
\tkzBannerTri[keys]{number}{title}
\end{codehigh}

\begin{demohigh}[language=latex/latex3,style/main=teal!25,style/code=teal!25]
\tkzBannerTri{01}{Title of document}
\end{demohigh}

Available \MontreCode{keys} are:

\begin{itemize}
	\item \MontreCode{height} (\MontreCode{2.5em} by default)
	\item \MontreCode{width} (\MontreCode{\textbackslash linewidth} by default)
	\item \MontreCode{blockwidth} (\MontreCode{2.75em} by default, but can be set to \MontreCode{auto})
	\item \MontreCode{coltxt} (\MontreCode{white} by default)
	\item \MontreCode{fonttxt}
	\item \MontreCode{swap} (\MontreCode{false} by default, for an other style )
	\item \MontreCode{maincolor} (\MontreCode{darkgray} by default)
	\item \MontreCode{collight} (\MontreCode{darkgray!25} by default)
	\item \MontreCode{colmedium} (\MontreCode{darkgray!50} by default)
	\item \MontreCode{coldark} (\MontreCode{darkgray} by default)
	\item \MontreCode{logo}
	\item \MontreCode{type}
	\item \MontreCode{dispblock} (\MontreCode{true} by default)
	\item \MontreCode{num} (\MontreCode{true}  by default)
	\item \MontreCode{customtype}
	\item \MontreCode{custommulti} (\MontreCode{false} by default)
\end{itemize}

\subsection{Examples}

\begin{demohigh}[language=latex/latex3,style/main=teal!25,style/code=teal!25]
\tkzBannerTri
  [maincolor=red,type=EXERCISES,blockwidth=auto,logo=\faAddressBook]
  {7}{My doc}
\end{demohigh}

\begin{demohigh}[language=latex/latex3,style/main=teal!25,style/code=teal!25]
\tkzBannerTri
  [maincolor=red,type=EXERCISES,blockwidth=5em,logo=\faAddressBook]
  {7}{My doc}
\end{demohigh}

\begin{demohigh}[language=latex/latex3,style/main=teal!25,style/code=teal!25]
\tkzBannerTri
  [maincolor=red,type=EXERCISES,blockwidth=auto,logo=\faAddressBook,swap]
  {07}{My doc}
\end{demohigh}

\begin{demohigh}[language=latex/latex3,style/main=teal!25,style/code=teal!25]
\tkzBannerTri
  [dispblock=false,maincolor=teal,logo=\faSchool]
  {}{My doc}
\end{demohigh}

\begin{demohigh}[language=latex/latex3,style/main=teal!25,style/code=teal!25]
\tkzBannerTri
  [maincolor=olive,customtype=TP,blockwidth=4em,logo=\faAddressBook,height=4em]
  {7}{My doc}
\end{demohigh}

\begin{demohigh}[language=latex/latex3,style/main=teal!25,style/code=teal!25]
\tkzBannerTriAlt
  [maincolor=violet,type=UE3.1,blockwidth=1.25cm,logo=\faGraduationCap,height=1.25cm]
  {TP}{My doc}
\end{demohigh}

\pagebreak

\section{Various commands}

\subsection{Difficulty levels with stars (fontawesome5)}

\begin{codehigh}[language=latex/latex3,style/main=teal!25,style/code=teal!25]
\DiffLevelStars[max level (3)]{level}
\end{codehigh}

\begin{demohigh}[language=latex/latex3,style/main=teal!25,style/code=teal!25]
\DiffLevelStars{0}\par
\DiffLevelStars{2.5}\par
\textcolor{teal}{\LARGE\DiffLevelStars[5]{4}}\par
\DiffLevelStars[5]{1.5}\par
\end{demohigh}

\subsection{Difficulty levels with stars (tikz)}

\begin{codehigh}[language=latex/latex3,style/main=teal!25,style/code=teal!25]
\tkzLevelStars[colframe=...,colback=...,offset=...,maxlevel=...,valign=...]{level}
\end{codehigh}

\begin{demohigh}[language=latex/latex3,style/main=teal!25,style/code=teal!25]
\tkzLevelStars{2.5}\par
{\LARGE We ty inline \tkzLevelStars{2.25} with score 2.25}\par
{\LARGE We ty inline \tkzLevelStars[valign=false]{1.75} with score 1.75}\par
\tkzLevelStars[colframe=red,colback=yellow,maxlevel=5]{3}
\end{demohigh}

\subsection{Flared arrow}

\begin{codehigh}[language=latex/latex3,style/main=teal!25,style/code=teal!25]
\tkzFlaredArrow[%
  color=...,           %color of arrow
  arrowsize=...,       %size (auto or H/W )
  bend=...,            %empty for straigth or left/... or right/...
  thickness=...,       %size for the beginning
  factor=...,          %factor for calculing size for ending
  arrowstyle=...,      %style (arrows.meta)
  move=...             %boolean for moving instead coordinates
  ]%
  {begin}{end or move}
\end{codehigh}

\begin{demohigh}[language=latex/latex3,style/main=teal!25,style/code=teal!25]
%arrow 0.5mm -> 1.25mm
\begin{tikzpicture}
\tkzFlaredArrow%
  [thickness=0.5mm,factor=2.5,bend=left/30,color=red,arrowstyle=Triangle]%
  {0,0}{5,1.5}
\end{tikzpicture}
\end{demohigh}

\begin{demohigh}[language=latex/latex3,style/main=teal!25,style/code=teal!25]
\begin{tikzpicture}
  \draw[thin,lightgray] (-3,-1) grid (5,5);
  \coordinate (A) at (0,0); \coordinate (B) at (4,1);
  \coordinate (C) at (1,1); \coordinate (D) at (5,4);
  \coordinate (E) at (0,1); \coordinate (F) at (0,5);
  \coordinate (G) at (-2,0);
  \tkzFlaredArrow[color=green,arrowstyle=Triangle]{A}{B}
  \tkzFlaredArrow[color=blue,bend=right/10]{D}{C}
  \tkzFlaredArrow%
    [color=red,bend=left/45,arrowstyle=Stealth,thickness=0.1mm,factor=10]%
    {-2,1}{0,4}
  \tkzFlaredArrow%
    [color=red,bend=right/45,thickness=0.1mm,factor=10,arrowstyle=Stealth]%
    {-2,1}{0,4}
  \tkzFlaredArrow[color=teal,move,bend=left/10]{-3,-1}{5,1}
  \end{tikzpicture}
\end{demohigh}

\subsection{Small markerbox}

\begin{codehigh}[language=latex/latex3,style/main=teal!25,style/code=teal!25]
\tbcmarker[color=...,width=...,font=...]{text}
\end{codehigh}

\begin{demohigh}[language=latex/latex3,style/main=teal!25,style/code=teal!25]
\tbcmarker{my text}
\end{demohigh}

\begin{demohigh}[language=latex/latex3,style/main=teal!25,style/code=teal!25]
\tbcmarker[color=olive,font=\normalfont\normalsize]{my text}
\end{demohigh}

%\subsection{Letters wirh pixelart style (experimental)}
%
%\begin{codehigh}[language=latex/latex3,style/main=teal!25,style/code=teal!25]
%%one letter only
%\PixlLetter%
%    [height=...,thick=...,color=...,gridcolor=...,
%    offseth=...,offsetv=...,gridafter=...,nospaceafter=...]
%    <tikzpicture options>
%    {letter}
%
%%several words
%\PixlLetters%
%    [height=...,thick=...,color=...,gridcolor=...,
%    offseth=...,offsetv=...,gridafter=...,nospaceafter=...]
%    {letters}
%\end{codehigh}
%
%\begin{codehigh}[language=latex/latex3,style/main=teal!25,style/code=teal!25]
%\PixlLetter{M}\PixlLetter{o}\PixlLetter{n}\PixlLetter{k}\PixlLetter{e}\PixlLetter{y}
%\end{codehigh}
%
%\PixlLetter{M}\PixlLetter{o}\PixlLetter{n}\PixlLetter{k}\PixlLetter{e}\PixlLetter{y}
%
%\begin{codehigh}[language=latex/latex3,style/main=teal!25,style/code=teal!25]
%\PixlLetters[height=2cm,color=red,gridafter,offsetv=1]{(Monkey)}
%\end{codehigh}
%
%\PixlLetters[height=2cm,color=red,gridafter,offsetv=1]{(Monkey)}
%
%\begin{codehigh}[language=latex/latex3,style/main=teal!25,style/code=teal!25]
%\PixlLetters[color=blue]{Monkeys roxxxx !?}
%\end{codehigh}
%
%\PixlLetters[color=blue]{Monkeys roxxxx !?}
%
%\begin{codehigh}[language=latex/latex3,style/main=teal!25,style/code=teal!25]
%\PixlLetters[color=red]{1+1=2}
%\end{codehigh}
%
%\PixlLetters[color=red]{1+1=2}

%\subsection{Wheel of skills / speedometer}
%
%\begin{codehigh}[language=latex/latex3,style/main=teal!25,style/code=teal!25]
%\WheelOfSkills[%
%    Radius=...,      %radius of wheel
%    Mark=...,        %absolute position of optional marker
%    Font=...,        %font of optional labels
%    SkillsList=...,  %list of optional skill labels
%    ]%
%    {number of skills or list of colors}
%\end{codehigh}
%
%\begin{demohigh}[language=latex/latex3,style/main=teal!25,style/code=teal!25]
%\WheelOfSkills[%
%    Radius=3cm,%
%    Mark=5.85,%
%    Font=\scriptsize\bfseries\ttfamily,%
%    SkillsList={Niv.1,Niv.2,Niv.3,Niv.4,Niv.5,Niv.6,Niv.7,Niv.8,Niv.9,Niv.10}]%
%    {10}%
%\end{demohigh}
%
%\begin{demohigh}[language=latex/latex3,style/main=teal!25,style/code=teal!25]
%\WheelOfSkills[%
%    Mark=1.5,%
%    Font=\scriptsize\bfseries\sffamily,%
%    SkillsList={LOW-MODERATE,NORMAL,HIGH,VERY HIGH,SEVERE,EXTREME,CATASTROPHIC}
%    ]%
%    {yellow!50,orange!50,red!50,blue!50,teal!50,purple!50,violet!50}%
%\end{demohigh}
%
%\begin{codehigh}[language=latex/latex3,style/main=teal!25,style/code=teal!25]
%\begin{SkillsWheel}[%
%    Radius=...,      %radius of wheel
%    Mark=...,        %absolute position of optional marker
%    Font=...,        %font of optional labels
%    SkillsList=...,  %list of optional skill labels
%    ]{number of skills or list of colors}
%    \PutIconsSkills[Pos=...,Scale=...]{list of icons}
%\end{SkillsWheel}
%\end{codehigh}
%
%\begin{demohigh}[language=latex/latex3,style/main=teal!25,style/code=teal!25]
%\begin{SkillsWheel}[Radius=5cm,Mark=5.85]%
%        {yellow!50,orange!50,red!50,blue!50,teal!50,purple!50,violet!50}
%    \PutIconsSkills[Echelle=3]%
%    {\faPython,\faAdjust,\faAngellist,\faAmbulance,\faAdjust,\faBabyCarriage,\faBlender}
%\end{SkillsWheel}
%\end{demohigh}
%
%\begin{codehigh}[language=latex/latex3,style/main=teal!25,style/code=teal!25]
%%inline version, with automatic dimensions
%\miniskillwheel[Colors=...,Mark=...]{nb of levels}
%
%%normal version
%\tkzspeedometer[Size=...,Mark=...,Colors=...]{nb levels}
%\end{codehigh}
%
%\begin{demohigh}[language=latex/latex3,style/main=teal!25,style/code=teal!25]
%%inline version, with automatic dimensions
%\scalebox{2.25}[2.25]{\sffamily Small inline \textit{skillwheel}
%\miniskillwheel[Colors=red/blue,Mark=4.33]{7} for testing.}
%
%%normal version
%\tkzspeedometer[Size=5cm,Mark=2.25,Colors=teal/magenta]{6}
%\end{demohigh}

\subsection{Annotate an image}

The idea is to provide a way of annotating an image, using an environment and a command which are linked to \TikZ.

\begin{codehigh}[language=latex/latex3,style/main=teal!25,style/code=teal!25]
\begin{imgannotate}[keys][includegraphics options]{imagefile with extension}
  \puttxtonimg[tikz node options]{coordinates}{txt}
  \puttxtonimg*[tikz node options]{coordinates within percentage}{txt}
\end{imgannotate}
%====keys
%clip=...     : boolean for clipping img
%node=...     : node name for reusing (remember picture)
%grid=...     : optionnal value for showing helping grid
%subgrid=...  : integer value for subgrid
%gridcolor=...: grid color
\end{codehigh}

\begin{demohigh}[language=latex/latex3,style/main=teal!25,style/code=teal!25]
%\usepackage[auto,outline]{contour}
\begin{imgannotate}[grid=1][height=4cm]{example-image.png}
  \puttxtonimg[scale=5,rotate=30]
    {1,1}{\contourlength{0.05em}\color{white}\contour{black}{$\pi$}}
  \puttxtonimg*[scale=1.5,rotate=-15]
    {0.66,0.75}{\contourlength{0.025em}\color{violet}\contour{yellow}{Pythagore}}
\end{imgannotate}
\end{demohigh}

\begin{demohigh}[language=latex/latex3,style/main=teal!25,style/code=teal!25]
\begin{imgannotate}[node=IMGTEST][height=5cm]<rotate=10>{example-image-golden.pdf}
  %tikz usual commands
  \draw (IMGTEST.center) node[above=5mm,font=\Huge\ttfamily\bfseries,text=blue] {TEST1};
  \draw (IMGTEST.north west) node[draw,thick,red,inner sep=0.5mm,below right=2.5mm,font=\LARGE\sffamily\bfseries,text=red] {TEST2};
\end{imgannotate}
\end{demohigh}

\subsection{Lengths}

\begin{codehigh}[language=latex/latex3,style/main=teal!25,style/code=teal!25]
\getwideststring[\macro]{elt1,elt2,...eltn}

\halignmakebox[align option]{elt}{list of elements}
\end{codehigh}

\begin{demohigh}[language=latex/latex3,style/main=teal!25,style/code=teal!25]
%widest string (\tmpwideststring by default)
\getwideststring{Exercise 1,Evaluation 2,Test n°3}\the\tmpwideststring
\end{demohigh}

\begin{demohigh}[language=latex/latex3,style/main=teal!25,style/code=teal!25]
%without
\sffamily\Large
Exercise 1 (10 points)\\
Evaluation 2 (8 points) \\
Test n°3 (4 points)
\end{demohigh}

\begin{demohigh}[language=latex/latex3,style/main=teal!25,style/code=teal!25]
%with
\sffamily\Large
\halignmakebox[l]{Exercise 1}{Exercise 1,Evaluation 2,Test n°3} 
(\halignmakebox[r]{10}{10,8,4} points)

\halignmakebox[l]{Evaluation 2}{Exercise 1,Evaluation 2,Test n°3} 
(\halignmakebox[r]{8}{10,8,4} points)

\halignmakebox[l]{Test n°3}{Exercise 1,Evaluation 2,Test n°3} 
(\halignmakebox[r]{4}{10,8,4} points)
\end{demohigh}

\begin{codehigh}[language=latex/latex3,style/main=teal!25,style/code=teal!25]
%width
\storewidthtolength[delta]{box}{\macro}
%height
\storeheighttolength[delta]{box}{\macro}
%totalheight
\storetotalheighttolength[delta]{box}{\macro}
%depth
\storedepthtolength[delta]{box}{\macro}
\end{codehigh}

\begin{demohigh}[language=latex/latex3,style/main=teal!25,style/code=teal!25]
\def\tmpbox{\large $1+\frac{1}{x}$}
%
\storewidthtolength{\tmpbox}{\mytmpboxwidth}\the\mytmpboxwidth

\storewidthtolength[10pt]{\tmpbox}{\mytmpboxwidthdelta}\the\mytmpboxwidthdelta

\storeheighttolength{\tmpbox}{\mytmpboxheight}\the\mytmpboxheight

\storetotalheighttolength{\tmpbox}{\mytmpboxtotheight}\the\mytmpboxtotheight

\storedepthtolength{\tmpbox}{\mytmpboxdepth}\the\mytmpboxdepth
\end{demohigh}

\begin{center}
\def\tmpbox{\large $1+\frac{1}{x}$}%
\storewidthtolength{\tmpbox}{\mytmpboxwidth}%
\storeheighttolength{\tmpbox}{\mytmpboxheight}%
\storetotalheighttolength{\tmpbox}{\mytmpboxtotheight}%
\storedepthtolength{\tmpbox}{\mytmpboxdepth}%
\begin{tikzpicture}[scale=5,transform shape]
	\node[draw=lightgray,inner sep=0pt,text=black] (BOX) {\large $1+\frac{1}{x}$};
	\begin{scope}
		\clip (BOX);
		\draw[lightgray] ([yshift=\mytmpboxdepth]BOX.south west) --++ ({\mytmpboxwidth+2\pgflinewidth},0);
	\end{scope}
	\draw[lightgray] ([yshift=\mytmpboxdepth]BOX.south west) --++ ({\mytmpboxwidth},0);
	\draw[<->,>=latex] ([yshift=-2pt]BOX.south west) --++ ({\mytmpboxwidth},0) node[midway,below,scale=0.2] {\the\mytmpboxwidth};
	\draw[<->,>=latex] ([xshift=-2pt]BOX.south west) --++ (0,\mytmpboxdepth) node[midway,left,scale=0.2] {\the\mytmpboxdepth};
	\draw[<->,>=latex] ([xshift=-2pt]BOX.north west) --++ (0,-\mytmpboxheight) node[midway,left,scale=0.2] {\the\mytmpboxheight};
	\draw[<->,>=latex] ([xshift=2pt]BOX.south east) --++ (0,\mytmpboxtotheight) node[midway,right,scale=0.2] {\the\mytmpboxtotheight};
\end{tikzpicture}
\end{center}

\begin{codehigh}[language=latex/latex2]
%starred version with box (tikz)
\fittexttobox(*){text}{width}{height}
\end{codehigh}

\begin{demohigh}[language=latex/latex2]
%with box
\fittexttobox*{PHONE}{2cm}{1cm}\\
\fittexttobox*{\bfseries\sffamily PHONE}{7cm}{1cm}\\
\fittexttobox*{PHONE}{3cm}{1cm}\\
\fittexttobox*{\ttfamily PHONE}{3cm}{1cm}\\
\fittexttobox*{PHONE}{2cm}{2cm}\\
\fittexttobox*{CONGRATULATIONS}{10cm}{3.5cm}\\
\fittexttobox*{CONGRATULATIONS}{14cm}{1.25cm}
\end{demohigh}

\begin{demohigh}[language=latex/latex2]
%w/o box
\fittexttobox{PHONE}{2cm}{1cm}\\
\fittexttobox{\bfseries\sffamily PHONE}{7cm}{1cm}\\
\fittexttobox{PHONE}{3cm}{1cm}\\
\fittexttobox{\ttfamily PHONE}{3cm}{1cm}\\
\fittexttobox{PHONE}{2cm}{2cm}\\
\fittexttobox{CONGRATULATIONS}{10cm}{3.5cm}\\
\fittexttobox{CONGRATULATIONS}{14cm}{1.25cm}
\end{demohigh}

\section{Small decorated boxes}

\subsection{Ornements box}

\begin{codehigh}[language=latex/latex2]
\begin{tcboxornaments}[keys]{options tcbox}
%inhalt
\end{tcboxornaments}

%---Available keys:
%v size deco        = height of deco
%h size deco        = width od deco
%linewidth deco     = thickness of deco
%color deco         = color of dedo
%alt deco           = boolean for inner deco
%h stretch deco     = coeff v -> h
%inner size deco    = corner size
\end{codehigh}

\begin{demohigh}[language=latex/latex2,style/main=teal!25,style/code=teal!25]
\begin{tcboxornaments}{}
\lipsum[1][1-5]
\end{tcboxornaments}
\end{demohigh}

\begin{demohigh}[language=latex/latex2,style/main=teal!25,style/code=teal!25]
\begin{tcboxornaments}
  [color deco=red!50!black,h size deco=3cm]
  {center,width=12cm,fontupper=\LARGE\sffamily,colupper=red!50!black}
CHAPTER 01: Logic
\end{tcboxornaments}
\end{demohigh}

\pagebreak

\subsection{Post 'forum' box}

\begin{codehigh}[language=latex/latex2]
\begin{tcforumpost}[keys]{options tcbox}
%inhalt
\end{tcforumpost}

%---Available keys:
%height-pseudo  = height of pseudo/alt-pseudo box
%deco-length    = width of triangular deco
%left-margin
%pseudo
%right          = boolean for right position
%swap-title     = boolean for switching lines in avatar box
%alt-pseudo
%fill-color
%rule-color
%avatar         = code avatar (emoji, graphics, txt...)
\end{codehigh}

\begin{demohigh}[language=latex/latex2,style/main=teal!25,style/code=teal!25]
\begin{tcforumpost}{}
\lipsum[1][1-5]
\end{tcforumpost}
\end{demohigh}

\begin{codehigh}[language=latex/latex3]
%txt pseudo styles (within l3e 'syntax')
\ExplSyntaxOn
\tikzset{
forum~style~post~font~pseudo/.style={%
	inner~sep=0pt,font=\ttfamily\bfseries,text=white},
forum~style~post~font~altpseudo/.style={%
	inner~sep=0pt,font=\ttfamily,text=\g_tcbox_forum_post_rule_color},
}
\ExplSyntaxOff
\end{codehigh}

\ExplSyntaxOn
\tikzset{
  forum~style~post~font~pseudo/.style={%
    inner~sep=0pt,font=\ttfamily\bfseries,text=white},
  forum~style~post~font~altpseudo/.style={%
    inner~sep=0pt,font=\ttfamily,text=\g_tcbox_forum_post_rule_color},
}
\ExplSyntaxOff

\begin{demohigh}[language=latex/latex2,style/main=teal!25,style/code=teal!25]
\begin{tcforumpost}%
  [%
    fill-color=pink!10,rule-color=magenta,%
    right,swap-title
  ]%
  {width=12cm}
\lipsum[1][1-5]
\end{tcforumpost}
\end{demohigh}

\end{document}